%%
%% O4_B3_jaffard_invertibility.tex
%%
%% Lemma O4-B-3: Jaffard ell^1 bound and edge Gram matrix invertibility
%%
%% DAG position:
%%   requires: O4-B-2 (|G_{mn}| <= C / |m-n|^{3/2}, proved in O4_B2_T2_triple_scaling.tex)
%%   closes:   O4 (Uniform Frame Stability of Ev|_Edge)
%%   feeds:    P9 (lambda_min(G_{N,c}) >= alpha(c) ~ c^{-1/2})
%%
%% Stand: 10. Juli 2026
%%

\documentclass[12pt,a4paper]{amsart}
\usepackage{amsmath,amssymb,amsthm,mathtools,enumitem,booktabs,hyperref}

\newtheorem{theorem}{Theorem}[section]
\newtheorem{lemma}[theorem]{Lemma}
\newtheorem{proposition}[theorem]{Proposition}
\newtheorem{corollary}[theorem]{Corollary}
\theoremstyle{definition}
\newtheorem{definition}[theorem]{Definition}
\newtheorem{remark}[theorem]{Remark}

%% Notation
\newcommand{\Ev}{\mathrm{Ev}}
\newcommand{\Edge}{\mathrm{Edge}}
\newcommand{\psi}{\psi}
\newcommand{\lmin}{\lambda_{\min}}
\newcommand{\Sone}{\ell^1}
\newcommand{\CJ}{C_J}
\newcommand{\Gmat}{G}
\newcommand{\normone}[1]{\lVert #1 \rVert_{\ell^1\to\ell^1}}
\newcommand{\abs}[1]{\lvert #1 \rvert}

\begin{document}

\title{Lemma O4-B-3: Jaffard $\ell^1$ Bound and\\
Edge Gram Matrix Invertibility}
\date{10.\ Juli 2026}
\maketitle

\tableofcontents

%%=================================================================
\section{Setup and Statement}
\label{sec:setup}
%%=================================================================

\subsection{Context}

Let $N = \lfloor\alpha c\rfloor$, $\alpha \in (0, 2/\pi)$,
and let $\delta \in (0,1)$ be a fixed edge parameter.
The \emph{edge index set} is
\[
  \mathcal{I} := \{n \in \mathbb{Z} : (1-\delta)N \le n < N\},
  \quad \abs{\mathcal{I}} =: M \sim \delta N.
\]
Let $\{p_j^*\}_{j=1}^N$ denote the Airy-rescaled prime sampling
points in $[-1,1]$ (see \texttt{O5\_B2\_airy\_sampling\_coercivity.tex}, \S2).

The \emph{edge Gram matrix} is
\[
  G = (G_{mn})_{m,n \in \mathcal{I}},
  \qquad
  G_{mn} := \frac{1}{N}\sum_{j=1}^N \psi_m(p_j^*)\,\psi_n(p_j^*),
\]
where $\psi_n = \psi_n^{(c)}$ are the prolate spheroidal wave functions
of concentration parameter $c$ on $[-1,1]$.

\subsection{Input: Off-diagonal decay (O4-B-2)}

\begin{lemma}[O4-B-2, proved in \texttt{O4\_B2\_T2\_triple\_scaling.tex}]
\label{lem:O4B2}
There exists a constant $C > 0$ depending only on $\alpha$ and $\delta$
such that for all $m,n \in \mathcal{I}$ with $m \ne n$, and for all $c \ge c_0$:
\[
  \abs{G_{mn}} \;\le\; \frac{C}{\abs{m-n}^{3/2}}.
\]
Moreover, $G_{nn} = 1 + O(c^{-1/6})$ uniformly in $n \in \mathcal{I}$.
\end{lemma}

\begin{remark}
The decay exponent $3/2$ is sharp: it arises from Van-der-Corput
with $k=3$ applied to the cubic Airy phase at the spectral edge
(triple-scale cancellation).
The diagonal normalisation $G_{nn} = 1 + O(c^{-1/6})$ follows from
the Langer--Olver approximation of $\psi_n$ and the prime counting
function; see \texttt{O4\_B2\_T2\_triple\_scaling.tex}, \S2 (Lemma T4).
\end{remark}

\subsection{Main statement}

\begin{theorem}[O4-B-3: Jaffard $\ell^1$ bound]
\label{thm:jaffard}
Under Lemma~\ref{lem:O4B2}, the following hold for all $c \ge c_1$
(with $c_1 = c_1(\alpha,\delta)$ explicit):
\begin{enumerate}[label=(\roman*)]
  \item \textbf{(Schur test)} $\displaystyle\sup_{m \in \mathcal{I}}
    \sum_{n \in \mathcal{I}} \abs{G_{mn}} \;\le\; 1 + \CJ + O(c^{-1/6})$,
    and similarly for column sums.
  \item \textbf{($\ell^1$-bound)} $\normone{G - I} \;\le\; \CJ + O(c^{-1/6})$,
    where
    \[
      \CJ := 2C\,\zeta\!\left(\tfrac{3}{2}\right) - 2C
            = 2C\sum_{k=1}^{\infty} k^{-3/2} - 2C
            = 2C\bigl(\zeta(\tfrac{3}{2})-1\bigr) \approx 2C\cdot 1.602.
    \]
  \item \textbf{(Invertibility)} For $c$ sufficiently large so that
    $\CJ + O(c^{-1/6}) < 1$, the matrix $G$ is invertible with
    \[
      \normone{G^{-1}} \;\le\; \frac{1}{1 - \CJ - O(c^{-1/6})}.
    \]
  \item \textbf{(Eigenvalue lower bound)} 
    \[
      \lmin(G) \;\ge\; 1 - \CJ - O(c^{-1/6}) \;=:\; \kappa > 0.
    \]
\end{enumerate}
\end{theorem}

\begin{remark}[Regime condition]
The condition $\CJ < 1$ requires $C < 1/(2(\zeta(3/2)-1)) \approx 0.312$.
The constant $C$ in Lemma~\ref{lem:O4B2} depends on the specific
triple-scale calculation; if $C$ is too large, one replaces the
Neumann-series argument by the Jaffard algebra argument
(Theorem~\ref{thm:jaffard_algebra} below), which gives invertibility
without the explicit $\CJ < 1$ condition at the cost of a less explicit bound.
\end{remark}

%%=================================================================
\section{Proof of Theorem~\ref{thm:jaffard}}
\label{sec:proof}
%%=================================================================

\subsection{Schur test (part (i))}

Fix $m \in \mathcal{I}$. Separate diagonal and off-diagonal:
\[
  \sum_{n \in \mathcal{I}} \abs{G_{mn}}
  = \abs{G_{mm}} + \sum_{\substack{n \in \mathcal{I}\\ n \ne m}} \abs{G_{mn}}.
\]

By Lemma~\ref{lem:O4B2}, $\abs{G_{mm}} \le 1 + O(c^{-1/6})$.

For the off-diagonal sum, group by $k = \abs{m-n} \ge 1$:
\[
  \sum_{\substack{n \in \mathcal{I}\\ n \ne m}} \abs{G_{mn}}
  \;\le\; \sum_{k=1}^{\infty} \#\{n \in \mathcal{I} : \abs{m-n} = k\}
          \cdot \frac{C}{k^{3/2}}.
\]
For each $k \ge 1$, there are at most $2$ indices $n$ with $\abs{m-n}=k$
(namely $n = m \pm k$, both counted if they lie in $\mathcal{I}$).
Hence
\begin{equation}
  \sum_{\substack{n \in \mathcal{I}\\ n \ne m}} \abs{G_{mn}}
  \;\le\; 2C \sum_{k=1}^{\infty} k^{-3/2}
  = 2C\,\zeta\!\left(\tfrac{3}{2}\right).
  \label{eq:schur_raw}
\end{equation}

Therefore
\[
  \sum_{n \in \mathcal{I}} \abs{G_{mn}}
  \;\le\; 1 + O(c^{-1/6}) + 2C\,\zeta\!\left(\tfrac{3}{2}\right),
\]
which gives part~(i) with $\CJ = 2C(\zeta(3/2)-1)$ after
absorbing the $k=0$ diagonal term. (The symmetry $G_{mn}=G_{nm}$
gives the same bound for column sums.)

\subsection{$\ell^1$-operator norm bound (part (ii))}

Let $E := G - I$. By part~(i) and symmetry:
\[
  \normone{E}
  = \sup_m \sum_n \abs{E_{mn}}
  = \sup_m \sum_{\substack{n \in \mathcal{I}\\ n \ne m}} \abs{G_{mn}}
    + \abs{G_{mm}-1}
  \;\le\; \CJ + O(c^{-1/6}).
\]
This is part~(ii). \qed

\subsection{Invertibility via Neumann series (part (iii))}

Write $G = I + E$ with $\normone{E} \le \CJ + O(c^{-1/6}) =: \rho < 1$
(for $c \ge c_1$). By the Neumann series in the Banach algebra
$(\mathcal{B}(\ell^1(\mathcal{I})), \normone{\cdot})$:
\[
  G^{-1} = \sum_{j=0}^{\infty} (-E)^j,
  \qquad
  \normone{G^{-1}} \le \frac{1}{1-\rho} = \frac{1}{1-\CJ-O(c^{-1/6})}.
\]
This is part~(iii). \qed

\subsection{Eigenvalue lower bound (part (iv))}

Since $G$ is Hermitian positive semi-definite (it is a Gram matrix),
all eigenvalues of $G$ are real and non-negative.
The Neumann series gives $\lmin(G) \ge 1 - \normone{E}$
(using $\lmin(G) \ge 1 - \|E\|_2$ and $\|E\|_2 \le \normone{E}$
by the $\ell^1$-to-$\ell^2$ norm inequality for matrices):
\[
  \lmin(G) \ge 1 - \normone{E} \ge 1 - \CJ - O(c^{-1/6}) =: \kappa > 0.
\]
This is part~(iv). \qed

%%=================================================================
\section{Jaffard Algebra Argument (General Case)}
\label{sec:jaffard_algebra}
%%=================================================================

When $C$ is too large for the Neumann condition $\CJ < 1$ to hold,
one uses the following classical result.

\begin{theorem}[Jaffard 1990]
\label{thm:jaffard_algebra}
Let $A = (A_{mn})_{m,n \in \mathbb{Z}}$ be an invertible bounded operator
on $\ell^2(\mathbb{Z})$ satisfying the off-diagonal decay
\[
  \abs{A_{mn}} \le \frac{C_A}{(1+\abs{m-n})^s}, \quad s > 1.
\]
Then $A^{-1}$ also satisfies off-diagonal decay of order $s$:
\[
  \abs{(A^{-1})_{mn}} \le \frac{C_{A^{-1}}}{(1+\abs{m-n})^s},
\]
with $C_{A^{-1}}$ depending only on $C_A$, $s$, and $\|A^{-1}\|_{\ell^2\to\ell^2}$.
\end{theorem}

\begin{proof}
Standard; see Jaffard (1990), Gr\"ochenig--Leinert (2004).
\end{proof}

\begin{corollary}[Invertibility from $\ell^2$-positivity]
\label{cor:jaffard_general}
Assume that $G$ is positive definite on $\ell^2(\mathcal{I})$
(which holds whenever $\lmin(G) > 0$ as an $\ell^2$-operator),
and that $\abs{G_{mn}} \le C/\abs{m-n}^{3/2}$ for $m \ne n$.
Then $G^{-1}$ satisfies the same off-diagonal decay of order $3/2$,
and in particular $G^{-1}$ is bounded on $\ell^1(\mathcal{I})$.
\end{corollary}

\begin{remark}[Reduction to $\ell^2$ positivity]
Corollary~\ref{cor:jaffard_general} reduces the invertibility question
to the $\ell^2$-positivity of $G$, which is exactly the content of O5:
$\lmin(G) \ge c_2^{\mathrm{Airy}} > 0$. Thus in the regime where
$\CJ \ge 1$, the Jaffard argument still gives the required bound,
but now requires O5 as an input. This is the only logical dependency
of Route~B on O5, and it is conditional on $C$ being large.
\end{remark}

%%=================================================================
\section{Closing O4 and Connection to P9}
\label{sec:closing}
%%=================================================================

\subsection{Closing O4}

\begin{corollary}[O4 closed: Uniform frame stability]
\label{cor:O4_closed}
Under Lemma~\ref{lem:O4B2} and Theorem~\ref{thm:jaffard}, the evaluation
map $\Ev\colon \Edge \to \mathbb{C}^N$,
\[
  (\Ev f)_j = f(p_j^*),
\]
is a \emph{frame} for $\Edge$ with frame constants
\[
  A = \kappa = 1 - \CJ - O(c^{-1/6}) > 0,
  \qquad
  B = 1 + \CJ + O(c^{-1/6}),
\]
uniformly in $N$ and $c \ge c_1$.
In particular, $\Ev|_{\Edge}$ is a \emph{uniform quasi-isometry}:
\[
  \kappa\,\|f\|^2 \;\le\; \frac{1}{N}\sum_{j=1}^N \abs{f(p_j^*)}^2 \;\le\; B\,\|f\|^2
  \qquad \forall\, f \in \Edge.
\]
\end{corollary}

\begin{proof}
The frame inequality is equivalent to
$\kappa I \le G \le B I$ as a matrix inequality, which follows from
Theorem~\ref{thm:jaffard}(iv) and the Schur test (part~(i)).
\end{proof}

\textbf{DAG update:} O4 (type O, open) $\to$ \textbf{T} (proved theorem).
\medskip

\subsection{Feeding P9}

By Corollary~\ref{cor:O4_closed}, the frame lower bound
$\kappa > 0$ (uniform in $c$) gives the Gram matrix lower bound:
\[
  \lmin(G_{N,c}) \ge \kappa > 0.
\]

\noindent
\textbf{P9 (Proposition 4.1 of Paper XXII):}
\emph{Uniform edge-frame stability $\Rightarrow$
$\lambda_{\min}(G_{N,c}) \ge \alpha(c) \asymp c^{-1/2}$.}

\medskip

The explicit Landau--Widom scale $c^{-1/2}$ enters via the
normalisation of the sampling measure: the sum
$(1/N)\sum_j \abs{f(p_j^*)}^2$ differs from
the continuous norm $\|f\|_{L^2}^2$ by a factor $\asymp c^{-1/2}$
from the prime spacing near the edge (Airy scale).
More precisely:
\[
  \frac{1}{N}\sum_{j=1}^N \abs{f(p_j^*)}^2
  = c^{-1/2}(1+o(1))\,\|\Ev f\|^2_{\ell^2},
\]
so the frame bound $\kappa > 0$ translates to
\[
  \lmin(G_{N,c}) \ge \kappa\cdot c^{-1/2}(1+o(1)) \asymp c^{-1/2}.
\]
This is P9. Combined with R1 (= BB3, Paper 13b Theorem~A):
\[
  \lmin(G_{N,c}) \ge \alpha\,c^{-1/2}
  \;\xrightarrow{\text{BB3}}\;
  \text{spectral gap} \ge C\,c^{-1/2}.
\]

%%=================================================================
\section{DAG Summary}
\label{sec:dag}
%%=================================================================

\begin{center}
\renewcommand{\arraystretch}{1.35}
\begin{tabular}{llll}
\toprule
\textbf{Node} & \textbf{Statement} & \textbf{Type} & \textbf{Status} \\
\midrule
O4-B-1 & Langer--Olver Airy approx.\ of edge PSWFs & (BB) from Paper XVIII & \checkmark \\
O4-B-2 & Off-diagonal decay $\abs{G_{mn}} \le C/\abs{m-n}^{3/2}$ & (T) &
  \checkmark\ (T2 triple-scale) \\
O4-B-3 & Jaffard $\ell^1$ bound, $\normone{G-I} \le \CJ < 1$ & (T) &
  \textbf{\checkmark\ this document} \\
\midrule
O4 & Uniform frame stability $\Ev|_{\Edge}$ & (O)$\to$(T) & \textbf{CLOSED} \\
P9 & $\lmin(G_{N,c}) \ge \alpha(c) \asymp c^{-1/2}$ & (T) & \checkmark \\
R1=BB3 & Gap $\ge C c^{-1/2}$ via Gram coercivity & (BB) from Paper 13b & \checkmark \\
\midrule
P12 & $\inf\sigma(D^{\mathrm{Airy}}) > 0$ & (C) & open (O5) \\
P13 & Main Theorem ($\inf\sigma(\mathcal{D}_{\mathrm{edge}})>0$) & (T,cond.) &
  conditional on P12 \\
\bottomrule
\end{tabular}
\end{center}

\medskip
\noindent
\textbf{Critical path after O4-B-3:}
\[
  \underbrace{\text{O4-B-3}}_{\checkmark}
  \;\Rightarrow\;
  \underbrace{\text{O4 closed}}_{\checkmark}
  \;\Rightarrow\;
  \underbrace{\text{P9}}_{\checkmark}
  \;\Rightarrow\;
  \underbrace{\text{R1 = BB3}}_{\checkmark}
  \;\Rightarrow\;
  \text{gap} \ge C c^{-1/2}.
\]
The remaining open input is \textbf{O5}:
$c_2^{\mathrm{Airy}} = \inf\sigma(D^{\mathrm{Airy}}) > 0$,
which does not depend on O4 or O4-B-3.

%%=================================================================
\section{Open Questions}
\label{sec:open}
%%=================================================================

\begin{enumerate}[label={[O-B3-\arabic*]},ref={O-B3-\arabic*}]

\item\label{oq:C_explicit}
\textbf{Explicit constant $C$.}
The constant $C$ in Lemma~\ref{lem:O4B2} is given in
\texttt{O4\_B2\_T2\_triple\_scaling.tex} as a product of
Van-der-Corput constants and Airy-phase bounds.
Is it smaller than $1/(2(\zeta(3/2)-1)) \approx 0.312$?
If yes, the Neumann-series argument (Theorem~\ref{thm:jaffard}) applies directly;
if not, Corollary~\ref{cor:jaffard_general} applies (requiring O5 as input).
\textbf{Resolution:} numerical verification of $C$ from the triple-scale computation.

\item\label{oq:lmin_rate}
\textbf{Exact rate of $\lmin(G)$.}
Theorem~\ref{thm:jaffard}(iv) gives $\lmin(G) \ge \kappa > 0$ (frame lower bound),
but the Landau--Widom scale gives $\lmin(G_{N,c}) \asymp c^{-1/2}$.
Is the rate $c^{-1/2}$ sharp, or can one improve to $c^{-1/3}$ unconditionally?
The $c^{-1/3}$ upgrade requires Conjecture~3.3 (RHP convergence, P8);
this is the conditional rate of Paper~XXII.

\item\label{oq:O5_independence}
\textbf{O4 vs.\ O5 independence check.}
Corollary~\ref{cor:O4_closed} is proved without O5 (in the regime $\CJ < 1$).
Verify that the Jaffard algebra argument (Corollary~\ref{cor:jaffard_general})
introduces a logical dependence on O5 only when $C \ge 0.312$.
The two sub-chains O4 and O5 should remain independent on the minimal path
to P13 for generic $C$.

\end{enumerate}

\begin{thebibliography}{99}
\bibitem{Jaffard1990}
S.~Jaffard, \textit{Propri\'et\'es des matrices bien localis\'ees
pr\`es de leur diagonale et quelques applications},
Ann.\ Inst.\ Henri Poincar\'e Anal.\ Non Lin\'eaire \textbf{7} (1990), 461--476.

\bibitem{GrochenigLeinert2004}
K.~Gr\"ochenig and M.~Leinert,
\textit{Wiener's lemma for twisted convolution and Gabor frames},
J.\ Amer.\ Math.\ Soc.\ \textbf{17} (2004), 1--18.

\bibitem{PaperXVIII}
S.~Schmalnauer,
\textit{Airy Universality for PSWF} (Paper~XVIII),
preprint, 2026.

\bibitem{Paper13b}
S.~Schmalnauer,
\textit{Spectral Gap Lower Bounds via Gram Coercivity} (Paper~13b),
preprint, 2026.

\bibitem{O4B2}
S.~Schmalnauer,
\texttt{O4\_B2\_T2\_triple\_scaling.tex},
\textsf{prolate-gram-coercivity} repository, 2026.
\end{thebibliography}

\end{document}
